\section{Race 1 --- trading the label directly}
\label{sec:race1}

\aidraft{%
The first and most direct use is the one the literature reports: hold the risky asset in the calm
state, de-risk in the stressed state.
We run the $K{=}2$ jump model on French daily total returns (trained from 1970, scored
1990-03--2026-05, $n{=}\overlapDays{}$ out-of-sample days) and evaluate the overlay against
buy-and-hold, against volatility targeting as the honest reactive incumbent, and against three
exposure-matched null distributions.%
}

\aidraft{%
The headline looks like the literature's: a \feeJMBH{}~bps/yr quadratic-utility fee over
buy-and-hold (Table~\ref{tab:exp1}).
It does not survive a single honest control.
The point estimate sits \emph{inside} every exposure-matched null band --- the persistent-placebo
95th percentile alone is $\placeboOne$, above the observed fee --- so random signals with matched
transition rates earn as much.
An exposure-matched static mix beats the overlay outright ($\feeStaticMix$), and at $\gamma{=}1$ the
fee turns to $\feeJMBHgone$.
Against the honest bar the verdict is unambiguous: $\text{fee}(\text{JM}-\text{VT}) = \feeJMVT$~bps/yr,
90\% CI $\feeJMVTci$, a significant loss to volatility targeting.
The frozen classification is Case~B: the beat-buy-and-hold result is an exposure/utility artifact,
and timing loses to the reactive incumbent.%
}

\begin{table}[t]
\centering
\small
\caption{Race 1 --- exposure overlay on equity: primary and controls. FKO quadratic-utility fee,
annualized bps/yr, $\gamma{=}10$ unless noted; 90\% CIs from the paired stationary bootstrap
(mean block 126d, $B{=}2000$). Source: \texttt{results/stage1.csv}.}
\label{tab:exp1}
\begin{tabular}{lr}
\toprule
Quantity & Value \\
\midrule
fee(JM $-$ B\&H), $\gamma{=}10$              & \feeJMBH~\feeJMBHci \\
fee(JM $-$ VT), $\gamma{=}10$ (honest bar)   & \feeJMVT~\feeJMVTci \\
fee(JM $-$ B\&H), $\gamma{=}1$               & \feeJMBHgone \\
fee(JM $-$ exposure-matched static mix)      & \feeStaticMix \\
same-close (delay $=1$) sensitivity          & $+480.1$ \\
break-even one-way cost                       & \breakeven~bps \\
C1 placebo 95th percentile ($n{=}100$)        & $+\placeboOne$ \\
C2 surrogate band max ($n{=}10$)              & $+1080.0$ \\
C3 iid band max ($n{=}10$)                    & $+805.3$ \\
label stability, train-start $+2$y / $-2$y    & \stabJM{} / \stabJM{} \\
era-half fees 1990--2007 / 2008--2026         & $+237.9$ / $+729.5$ \\
switches/yr                                    & \switchesYr{} \\
frozen classification                         & Case B \\
\bottomrule
\end{tabular}
\end{table}

